\documentclass[a4paper,12pt]{article}

%\documentclass[a4paper,12pt, draft]{article}
\documentclass[a4paper,12pt]{article}

%%% Работа с русским языком
\usepackage{cmap}					% поиск в PDF
\usepackage{mathtext} 				% русские буквы в формулах
\usepackage[T2A]{fontenc}			% кодировка
\usepackage[utf8]{inputenc}			% кодировка исходного текста
\usepackage[english,russian]{babel}	% локализация и переносы
\usepackage{indentfirst}            % красная строка в первом абзаце
\frenchspacing                      % равные пробелы между словами и предложениями

%%% Дополнительная работа с математикой
\usepackage{amsmath,amsfonts,amssymb,amsthm,mathtools} % пакеты AMS
\usepackage{icomma}                                    % "Умная" запятая

%%% Свои символы и команды
\usepackage{centernot} % центрированное зачеркивание символа
\usepackage{stmaryrd}  % некоторые спецсимволы

\renewcommand{\epsilon}{\ensuremath{\varepsilon}}
\renewcommand{\phi}{\ensuremath{\varphi}}
\renewcommand{\kappa}{\ensuremath{\varkappa}}
\renewcommand{\le}{\ensuremath{\leqslant}}
\renewcommand{\leq}{\ensuremath{\leqslant}}
\renewcommand{\ge}{\ensuremath{\geqslant}}
\renewcommand{\geq}{\ensuremath{\geqslant}}
\renewcommand{\emptyset}{\ensuremath{\varnothing}}

\DeclareMathOperator{\sgn}{sgn}
\DeclareMathOperator{\ke}{Ker}
\DeclareMathOperator{\im}{Im}
\DeclareMathOperator{\re}{Re}

\newcommand{\N}{\mathbb{N}}
\newcommand{\Z}{\mathbb{Z}}
\newcommand{\Q}{\mathbb{Q}}
\newcommand{\R}{\mathbb{R}}
\newcommand{\Cm}{\mathbb{C}}
\newcommand{\F}{\mathbb{F}}
\newcommand{\id}{\mathrm{id}}
\newcommand{\imp}[2]{
	(#1\,\,$\ra$\,\,#2)\,\,
}
\newcommand{\System}[1]{
	\left\{\begin{aligned}#1\end{aligned}\right.
}
\newcommand{\Root}[2]{
	\left\{\!\sqrt[#1]{#2}\right\}
}
\newcommand{\RR}{\R}
\newcommand{\NN}{\N}
\renewcommand{\subseteq}{\subset}
\newcommand{\sub}{\subset}
\newcommand{\sconstr}{\;\vert\;}
\newcommand{\thus}{\implies}

\newcommand{\defeq}{\vcentcolon= }
\newcommand{\defev}{\stackrel{\Delta}{\Longleftrightarrow}}
\newcommand{\deriv}[3][1]{%
	\ifthenelse{#1>1}{%
		\frac{\dlta^{#1} {#2}}{\dlta {#3}^{#1}}
	}{%
		\frac{\dlta {#2}}{\dlta {#3}}
	}%
}

\renewcommand\labelitemi{$\triangleright$}

\let\bs\backslash
\let\lra\Leftrightarrow
\let\ra\Rightarrow
\let\la\Leftarrow
\let\emb\hookrightarrow

%%% Перенос знаков в формулах (по Львовскому)
\newcommand{\hm}[1]{#1\nobreak\discretionary{}{\hbox{$\mathsurround=0pt #1$}}{}}

%%% Работа с картинками
\usepackage{graphicx}    % Для вставки рисунков
\setlength\fboxsep{3pt}  % Отступ рамки \fbox{} от рисунка
\setlength\fboxrule{1pt} % Толщина линий рамки \fbox{}
\usepackage{wrapfig}     % Обтекание рисунков текстом

%%% Работа с таблицами
\usepackage{array,tabularx,tabulary,booktabs} % Дополнительная работа с таблицами
\usepackage{longtable}                        % Длинные таблицы
\usepackage{multirow}                         % Слияние строк в таблице

%%% Теоремы
\theoremstyle{plain}
\newtheorem{theorem}{Теорема}[section]
\newtheorem{lemma}{Лемма}[section]
\newtheorem{proposition}{Утверждение}[section]
\newtheorem*{exercise}{Упражнение}
\newtheorem*{problem}{Задача}

\theoremstyle{definition}
\newtheorem{definition}{Определение}[section]
\newtheorem*{corollary}{Следствие}
\newtheorem*{note}{Замечание}
\newtheorem*{reminder}{Напоминание}
\newtheorem*{example}{Пример}

\theoremstyle{remark}
\newtheorem*{solution}{Решение}

%%% Оформление страницы
\usepackage{extsizes}     % Возможность сделать 14-й шрифт
\usepackage{geometry}     % Простой способ задавать поля
\usepackage{setspace}     % Интерлиньяж
\usepackage{enumitem}     % Настройка окружений itemize и enumerate
\setlist{leftmargin=25pt} % Отступы в itemize и enumerate

\geometry{top=25mm}    % Поля сверху страницы
\geometry{bottom=30mm} % Поля снизу страницы
\geometry{left=20mm}   % Поля слева страницы
\geometry{right=20mm}  % Поля справа страницы

\setlength\parindent{15pt}        % Устанавливает длину красной строки 15pt
\linespread{1.3}                  % Коэффициент межстрочного интервала
%\setlength{\parskip}{0.5em}      % Вертикальный интервал между абзацами
%\setcounter{secnumdepth}{0}      % Отключение нумерации разделов
%\setcounter{section}{-1}         % Нумерация секций с нуля
\usepackage{multicol}			  % Для текста в нескольких колонках
\usepackage{soulutf8}             % Модификаторы начертания
\mathtoolsset{showonlyrefs=true} % показывать номера формул только у тех, у которых есть ссылки по eqref
%%% Содержаниие
\usepackage{tocloft}
\tocloftpagestyle{main}
%\setlength{\cftsecnumwidth}{2.3em}
%\renewcommand{\cftsecdotsep}{1}
%\renewcommand{\cftsecpresnum}{\hfill}
%\renewcommand{\cftsecaftersnum}{\quad}

%%% Шаблонная информация для титульного листа
\newcommand{\CourseName}{Математический анализ}
\newcommand{\FullCourseNameFirstPart}{\so{Введение в математический анализ}}
\newcommand{\SemesterNumber}{I}
\newcommand{\LecturerInitials}{Гусев Николай Анатольевич}
\newcommand{\CourseDate}{осень 2022}
\newcommand{\AuthorInitials}{Копанов Антон}
\newcommand{\AuthorInitialssecond}{Клищ Данил}
\newcommand{\VKLink}{https://vk.com/notnako}
\newcommand{\GithubLink}{https://github.com/MIPT-Group/Lectures_Tex_Club}
\newcommand{\VKLinksecond}{https://vk.com/dan.klishch}

%%% Колонтитулы
\usepackage{titleps}
\newpagestyle{main}{
	\setheadrule{0.4pt}
	\sethead{\CourseName}{}{\hyperlink{intro}{\;Назад к содержанию}}
	\setfootrule{0.4pt}                       
	\setfoot{ФПМИ МФТИ, \CourseDate}{}{\thepage} 
}
\pagestyle{main}  

%%% Нумерация уравнений
\makeatletter
\def\eqref{\@ifstar\@eqref\@@eqref}
\def\@eqref#1{\textup{\tagform@{\ref*{#1}}}}
\def\@@eqref#1{\textup{\tagform@{\ref{#1}}}}
\makeatother                      % \eqref* без гиперссылки
\numberwithin{equation}{section}  % Нумерация вида (номер_секции).(номер_уравнения)
\mathtoolsset{showonlyrefs= true} % Номера только у формул с \eqref{} в тексте.

%%% Гиперссылки
\usepackage{hyperref}
\usepackage[usenames,dvipsnames,svgnames,table,rgb]{xcolor}
\hypersetup{
	unicode=true,            % русские буквы в раздела PDF
	colorlinks=true,       	 % Цветные ссылки вместо ссылок в рамках
	linkcolor=black!15!blue, % Внутренние ссылки
	citecolor=green,         % Ссылки на библиографию
	filecolor=magenta,       % Ссылки на файлы
	urlcolor=NavyBlue,       % Ссылки на URL
}

%%% Графика
\usepackage{tikz}        % Графический пакет tikz
\usepackage{tikz-cd}     % Коммутативные диаграммы
\usepackage{tkz-euclide} % Геометрия
\usepackage{stackengine} % Многострочные тексты в картинках
\usetikzlibrary{angles, babel, quotes}

%%% Всю шаблонную информацию можно менять тут
\newcommand{\FullCourseNameFirstPart}{\so{ТЕОРИЯ И РЕАЛИЗАЦИЯ}}
\newcommand{\FullCourseNameSecondPart}{\so{ЯЗЫКОВ ПРОГРАММИРОВАНИЯ}}
\newcommand{\SchoolName}{ФПМИ}
\newcommand{\TrackName}{МФ}
\newcommand{\SemesterNumber}{III}
\newcommand{\LecturerInitials}{Рубцов Александр Александрович}
\newcommand{\AutherInitials}{Моложавенко Александр}
\newcommand{\VKLink}{https://vk.com/}
\newcommand{\OverleafLink}{https://www.overleaf.com/}
\newcommand{\GithubLink}{https://github.com}
\newcommand{\CourseName}{Краткое Название} %  Используется в преамбуле для создания названия предмета в верхнем контитуле   
\newcommand{\CourseDate}{Осень 2020}           %  Используется в преамбуле для создания даты в нижнем контитуле и в title_page

%\includeonly{lectures/lect05,lectures/lect06}  % Чтобы скомпилировать только часть лекций

\begin{document}
%%% Всю шаблонную информацию можно менять тут
\newcommand{\FullCourseNameFirstPart}{\so{ТЕОРИЯ ГРУПП}}
\newcommand{\SemesterNumber}{III}
\newcommand{\LecturerInitials}{Богданов Илья Игоревич}
\newcommand{\AutherInitials}{Даниил Дрябин}
\newcommand{\VKLink}{https://vk.com/ax_equals_b}
\newcommand{\GithubLink}{https://github.com/daniild71r/lectures_tex_club}

\begin{titlepage}
	\clearpage\thispagestyle{empty}
	\centering
	
	%\textit{Федеральное государственное автономное учреждение \\
	%	высшего профессионального образования}
	%\vspace{0.5ex}
	
	%\textbf{Московский Физико-Технический Институт \\ КЛУБ ТЕХА ЛЕКЦИЙ}
	\textbf{Московский физико-технический институт \\ Физтех-школа прикладной математики и информатики}
	\vspace{20ex}
	\vspace{13ex}
	
	{\textbf{\FullCourseNameFirstPart}}
	
	\SemesterNumber\ СЕМЕСТР  
	\vspace{1ex}
	
	Лектор: \textit{\LecturerInitials}
	
	\begin{figure}[!ht]
		\centering
		\includegraphics[width=0.4\textwidth]{images/logo_ltc.png}
	\end{figure}

	\begin{flushright}
		\noindent
		Автор 2020: \href{\VKLink}{\textit{\AutherInitials}}
		\\
		Дополнил 2022: \textit{Даниил Максимов}
		\\
		\href{\GithubLink}{\textit{Проект на Github}}
	\end{flushright}
	
	\vfill
	\CourseDate
	\pagebreak
\end{titlepage}

\newpage
\hypertarget{intro}{}
\tableofcontents

\newpage
\section{Unix системы}

\subsection{Компоненты операционной системы}

\begin{Def}
	\underline{Ядро} --- программа, которая запускается самой первой. Имеет привилегированное положение на процессоре и может взаимодействовать с <<железом>> напрямую.
\end{Def}

\begin{Def}
	\underline{Базовые библиотеки}: зачастую входят в состав операционной системы, потому что привязаны к определенному ядру (версии ядра), чтобы максимально использовать его функциональность. 
\end{Def}

\begin{Def}
\underline{Служебные сервисы} --- ведение логов, обработка сетевых соединения, и т.д.
\end{Def}

\begin{Def}
\underline{Минимальная пользовательская оболочка} --- чтобы пользователь мог взаимодействовать с системой.
\end{Def}

\subsection{Процессы}

\begin{Def}
\underline{Процесс} --- любой экземпляр программы, работающей в операционной системе.
\end{Def}

Примеры:
\begin{itemize}
	\item Приложения
	\item Фоновые сервисы (демоны)
	\item Команды в терминале
\end{itemize}

Главные команды для мониторинга процессов:
\begin{itemize}
	\item ps -A -- список процессов
	\item pstree -- иерархия процессов
	\item top
\end{itemize}

\subsection{UNIX программы}

Программы:
\begin{itemize}
	\item \textbf{Бинарные исполняемые файлы.}
	\item \textbf{Текстовый файл, начинающийся с \#!.} После \#! указывается путь к интерпретатору (например, \#!/bin/bash). Это могут быть скрипты на Python/Perl/Escript или Shell скрипт.
\end{itemize}

Название файла и его расширение не имеет значения. Любой файл, имеющий атрибут 
<<исполняемый>> является программой.

\subsection{UNIX пользователи}

Единственные привилегированный пользователь: root (UID = 0).

Непривилегерованные пользователи:
\begin{itemize}
	\item Реальные пользователи, которые могут войти в систему (UID $\ge$ 1000)
	\item Фейковые пользователи, привязанные к сервисам (UID $<$ 1000).
\end{itemize}

\textbf{Зачем сервисам нужны фейковые пользователи?}
Допустим, у вас есть веб-сервер и сервер баз данных. Теоретически какой-то пользователь
может подключиться к веб-серверу и найти там уязвимость. Для того, чтобы 
минимизировать возможный ущерб от уязвимости в одном из процессов, процессы запускаются
под разными пользователями и не имеют права общаться друг с другом.

Непривилегерованные пользователи могут временно получить 
дополнительные привилегии:

\begin{itemize}
	\item su - запускает терминал от имени root
	\item sudo - запускает любую программу от имени root
	\item Запустить программу с атрибутом SUID и владельцем root. SUID атрибут означает, что при запуске программа
	исполняется от того пользователя, кто является владельцем файла. Кстати, так и 
	работает команда sudo, владелец sudo - root, а также у sudo проставлен атрибут SUID.
	\item Запустить программу с <<capabilities>> флагами. Эти флаги позволяют тонко настроить, что можно делать программе, а что - нельзя.
\end{itemize}

\subsection{UNIX межпроцессорное взаимодействие}

Процессы работают изолированно друг от друга. Единственный способ взаимодействовать друг с другом --- использовать методы межпроцессорного взаимодействия ядра (Inter-Process Communication, IPC):

\begin{itemize}
	\item {Сигналы} для коротких сообщений
	\item {Каналы и сокеты} для последовательных данных
	\item {Разделяемая память}, в которой хранятся большие куски данных
\end{itemize}

\subsection{Процесс запуска}

\begin{enumerate}
	\item {Загрузчик} с диска или UEFI материнской платы выполняется в привилегированном режиме
	\item Загрузчик запускает {ядро}, тоже в привилегированном режиме.
	\item Ядро запускает первый непривелигерованный процесс: {init или systemd}.
\end{enumerate}

\subsection{Сервисы (``Демоны'')}

\begin{Def}
	\underline{Сервисы} --- это специальные системные службы, которые работают в фоновом режиме.
\end{Def}

\newpage

Минимальный сет во многих системах:

\begin{itemize}
	\item \textbf{dhclient} -- держит активным полученный IP адрес в сети
	\item \textbf{getty} -- переключение между графическим/консольным входом в систему
	\item \textbf{sshd} -- подключение по ssh
	\item \textbf{ntpd или chronyd} -- синхронизация времени
	\item \textbf{syslogd} -- системный лог
\end{itemize}

\subsection{\#!/bin/sh}

В терминале вы работаете в интерпретаторе \texttt{shell}. \texttt{shell} предоставляет возможность выполнять
команды и последовательности команд.

Бывают разные интерпретаторы \texttt{shell}:
\begin{itemize}
	\item самодостаточный \texttt{shell} program: \textbf{FreeBSD}
	\item основной в \texttt{Linux}: \textbf{bash}
	\item debian: \textbf{dash}
	\item alpine: \textbf{busybox}
	\item maxOS X: \textbf{zsh}
\end{itemize}

\subsection{Системный демон (systemd)}
\begin{Def}
Процесс, который управляет всеми остальными процессами. При этом сам является обычным процессом. \texttt{Systemd} пришел на замену INIT-процессу.

\end{Def}

Чем systemd лучше?
\begin{itemize}
	\item Устранение интерпретации \texttt{shell} скриптов. 
	\item Запуск сервисов параллельно
\end{itemize}

\subsubsection{Концепты systemd}

\begin{itemize}
	\item Сервисы
	\item Цели:
		\begin{itemize}
			\item \underline{Стандартные цели}. Например, уровни \texttt{init} в System-V:
				\begin{itemize}
					\item poweroff.target (init 0)
					\item rescue.target (init 1)
					\item multi-user.target (init 3)
					\item graphical.target (init 5)
					\item reboot.target (init 6)
					\item sleep.target и hibertate.target (не присутствуют в System-V)
				\end{itemize}
			\item \underline{Специальные цели}
		\end{itemize}
\end{itemize}

\subsection{Фоновые задачи}

\begin{itemize}
	\item \textbf{Нет привязанного к задаче терминала.} Для того, чтобы сигнал \texttt{SIGHUP}
	не останавливал исполнение задачи используется команда \texttt{nohup}.
	\item \textbf{Пишут в лог.} Иногда пишут данные в какую-то старинную 
	локальную почтовую систему.
	\item \textbf{Запускается в единственном экземпляре.} Обычно это достигается
	с помощью file lock.
\end{itemize}

\subsection{Дистрибутивы \texttt{Linux}}

Ядро называется Linux. Базовая система --- минимальная юзабельная система.

\underline{Дистрибутив}:
\begin{itemize}
	\item Ядро
	\item Базовая система
	\item GUI оболочка
	\item Программное обеспечение 
\end{itemize}

\subsubsection{Какие есть дистрибутивы?} 

Общего назначения:
\begin{itemize}
	\item OpenSUSE
	\item Fedora
	\item Debian
	\item Ubuntu
\end{itemize}

Специфические дистрибутивы:
\begin{itemize}
	\item Alpine Linux --- очень легковесный, поэтому его
	хорошо ставить на виртуальную машину.
	\item Kali Linux --- используется для тестов на безопасность.
\end{itemize}


\newpage
\begin{flushright}
    \textit{Лекция 2 (от 08.09)}
\end{flushright}
\section{$\S 3.$ Дифференцирование}
Пусть $z_0 \in \CC$, $B_r(z_0)$ и $f: B_r(z_0) \mapsto \CC$
\Def
Пусть $f: B_r(z_0) \mapsto \CC$. Если
\begin{equation}\label{(3.1)}
    \exists \dst \lim_{z \mapsto z_0} \frac{f(z)-f(z_0)}{z-z_0}
\end{equation}
то он называется \textbf{производной функции $f$ в точке $z_0$} (записывается
$f'(z_0)$).
\\
Можем расписать предел:
\begin{equation}\label{(3.2)}
    \forall \varepsilon > 0 \ \exists \delta(\varepsilon)>0: \ z \in B_\delta(z_0), \ \left| \frac{f(z) - f(z_0)}{z-z_0} - f'(z_0) \right| < \varepsilon
\end{equation}
\begin{equation}\label{(3.3)}
    f(z) = f(z_0) + f'(z_0)(z-z_0) + \alpha(z-z_0)
\end{equation}
\begin{equation}\label{(3.4)}
    \lim_{z \to z_0}\frac{\alpha(z-z_0)}{z-z_0} = 0
\end{equation}
\begin{align*}
  & f(z_0) = u(x,y) + iv(x,y)
\end{align*}
\Def
Говорят, что $f: B_r(z_0) \mapsto \CC$ \textbf{дифференцируема в точке $z_0$},
если
\begin{equation}\label{(3.5)}
    \exists A \in \CC: \ f(z) = f(z_0)+A(z-z_0)+o(z-z_0)
\end{equation}
\lemma
$f: B_r(z_0) \mapsto \CC$ дифференцируема в $z_0$ $\Leftrightarrow$ $\exists
f'(z_0)$, $A = f'(z_0)$
\pr
Очевидно.
\\
Обозначим: $\Delta x = x-x_0$, $\Delta y = y - y_0$, $\Delta z = z-z_0 = \Delta
x + i \Delta y$, $\Delta u = u(x,y)-u(x_0,y_0)$, $\Delta v = v(x,y)-v(x_0,y_0)$,
$\Delta f = f(x,y)-f(x_0,y_0) = \Delta u + i \Delta v$.
\theorem
$f: B_r(z_0) \mapsto \CC$ дифференцируема в $z_0$ т огда и только тогда, когда
\begin{itemize}
    \item $u(x,y)$, $v(x,y)$ дифференцируемы в $(x_0,y_0)$ (в $\RR^2$)
    \item выполняется условие Коши-Римана (УКР):
    \begin{equation}\label{(3.6)}
        \begin{cases}
            \dst \frac{\partial u}{\partial x} = \dst \frac{\partial v}{\partial y} \\
            \dst \frac{\partial v}{\partial x} = - \dst \frac{\partial u}{\partial y}
        \end{cases}
        % \left\{ \begin{matrix}
        %         \dst \frac{\partial u}{\partial x} = \dst \frac{\partial v}{\partial y} \\
        %         \dst \frac{\partial v}{\partial x} = - \dst \frac{\partial u}{\partial y}
        %     \end{matrix} \right.
    \end{equation}
\end{itemize}
При этом
\begin{equation}\label{(3.7)}
    f'(z_0) = \frac{\partial u}{\partial x}(x_0,y_0) + i\frac{\partial v}{\partial x}(x_0,y_0) = \frac{\partial v}{\partial y}(x_0,y_0) - i\frac{\partial u}{\partial y}(x_0,y_0)
\end{equation}
\pr
Докажем в обе стороны.
\begin{itemize}
    \item Необходимость.
    \begin{align*}
      & \exists f'(z_0) = a+ib = A \in \CC
    \end{align*}
    Значит, по определению дифференцируемости
    \begin{align*}
      & \Delta f = A \Delta z + \alpha(\Delta z); \ \alpha(\Delta z) = \alpha_1(\Delta x, \Delta y) + i\alpha_2(\Delta x, \Delta y)
    \end{align*}
    \begin{align*}
      & \left\{ \begin{matrix}
              \Delta u = a \Delta x - b \Delta y + \alpha_1(\Delta x, \Delta y) \\
              \Delta v = b \Delta x + a \Delta y + \alpha_2(\Delta x, \Delta y)
          \end{matrix} \right.
    \end{align*}
    \begin{align*}
      & \left| \alpha_1 \right| \leq \left| \alpha \right|, \ \left| \alpha_2 \right| \leq \left| \alpha \right|
    \end{align*}
    \begin{align*}
      & \lim_{(\Delta x; \Delta y) \to (0;0)}\left( \frac{\left| \alpha_1(\Delta x, \Delta y) \right|}{\sqrt{\left( \Delta x \right)^2 + \left( \Delta y \right)^2}} \right) \leq \lim_{\Delta z \to 0}\left( \frac{\left| \alpha(\Delta z) \right|}{\left| \Delta z \right|} \right) = 0
    \end{align*}
    Значит, $u$ дифференцируема, причем $\dst \frac{\partial u}{\partial x} =
    a$, $\dst \frac{\partial u}{\partial y} = -b$.
    \\
    Аналогично для $v$, причем $\dst \frac{\partial v}{\partial x} = b$, $\dst
    \frac{\partial v}{\partial y} = a$.
    \\
    Видим выполнимость УКР.
    \\
    Выполнимость \eqref{(3.7)} очевидна.
    \item Достаточность.
    \\
    Пусть $u$, $v$ дифференцируемы в $(x_0,y_0)$ и выполняется УКР. Тогда
    \begin{align*}
      & \Delta f = \Delta u + i \Delta v = \frac{\partial u}{\partial x} \Delta x + \frac{\partial u}{\partial y} \Delta y + \alpha_1(\Delta x, \Delta y) + i \left( \frac{\partial v}{\partial x} \Delta x + \frac{\partial v}{\partial y} \Delta y + \alpha_2(\Delta x, \Delta y)\right) = \\
      & \frac{\partial u}{\partial x} \Delta x - \frac{\partial v}{\partial x} \Delta y + \alpha_1(\Delta x, \Delta y) + i \left( \frac{\partial v}{\partial x} \Delta x + \frac{\partial u}{\partial x} \Delta y + \alpha_2(\Delta x, \Delta y)\right) = \left( \frac{\partial u}{\partial x} + i \frac{\partial v}{\partial x}\right)\cdot \\
      & \cdot \left( \Delta x + i \Delta y \right)+ \alpha_1(\Delta x, \Delta y) + i \alpha_2(\Delta x, \Delta y)
    \end{align*}
    Значит,
    \begin{align*}
      & \exists f'(z_0) = \frac{\partial u}{\partial x} + i \frac{\partial v}{\partial x}
    \end{align*}
\end{itemize}
\Example
\begin{align*}
  & w = z^2
\end{align*}
\begin{itemize}
    \item[I] Найдем производную по определению.
    \begin{align*}
      & \frac{\Delta w}{\Delta z} = \frac{\left( z_0-\Delta z \right)^2 - z_0^2}{\Delta z} = \frac{2z_0\Delta z + \left( \Delta z \right)^2}{\Delta z} = 2z_0 + \Delta z \us{\Delta z \to 0}{\to} 2z_0
    \end{align*}
    \item[II] Найдем производную с использованием теоремы.
    \begin{align*}
      & z^2 = x^2-y^2+2ixy \Rightarrow u = x^2-y^2, \ v = 2xy
    \end{align*}
    \begin{align*}
      & u_x = 2x = v_y; \ u_y = -2y = - v_x
    \end{align*}
    Видим выполнимость УКР, поэтому
    \begin{align*}
      & w' = u_x+iv_x = 2x + 2iy = 2z
    \end{align*}
\end{itemize}
\Example
\begin{align*}
  & w = \left| z \right|^2
\end{align*}
\begin{align*}
  & \left| z \right|^2 = x^2+y^2 \Rightarrow u = x^2+y^2, \ v = 0
\end{align*}
\begin{align*}
  & u_x = 2x \neq 0 = v_y; \ u_y = 2y \neq 0 = - v_x
\end{align*}
(за исключением точки $(0;0)$). Видим выполнение УКР только в точке $(0;0)$, а
значит, и функция дифференцируема только в этой точке (с производной, равной
$0$).
\Example
\begin{align*}
  & w = \ol{z}
\end{align*}
\begin{align*}
  & \ol{z} = x - iy \Rightarrow u = x, \ v = -y
\end{align*}
\begin{align*}
  & u_x = 1 \neq -1 = v_y; \ u_y = 0 = - v_x
\end{align*}
Видим, что УКР не выполняется ни в одной точке $\CC$, а значит, функция нигде не
дифференцируема.
\prop
Если $f(z)$, $g(z)$ дифференцируемы, то:
\begin{itemize}
    \item дифференцируема и сумма:
    \begin{align*}
      & \left( f+g \right)' = f'+g'
    \end{align*}
    \item дифференцируемо и произведение:
    \begin{align*}
      & \left( fg \right)' = f'g+fg'
    \end{align*}
    \item дифференцируемо и частное (когда $g(z) \neq 0$):
    \begin{align*}
      & \left( \frac{f}{g} \right)' = \frac{f'g+fg'}{g^2}
    \end{align*}
\end{itemize}
\pr
Доказательства легко провести через УКР с использованием действительного
анализа.
\section{$\S 4.$ Регулярные и гармонические функции}
\example
$B_r(z_0)$~--- односвязная область как в $\CC$, так и в $\CCC$.
\example
$K = \left\{ z: 0 < r_1 < \left| z \right| < r_2 < \infty \right\}$~--- не
односвязная область ни в $\CC$, ни в $\CCC$.
\example
$B_r(\infty)$~--- односвязная область в $\CCC$, но не в $\CC$.
\Def
Если $u:G \mapsto \RR^1$, $G \subseteq \RR^2$~--- область, $u \in C^2(G)$,
$\Delta u = 0$, где $\Delta = \dst \frac{\partial^2}{\partial x^2} + \dst
\frac{\partial^2}{\partial y^2}$, то $u$ называется \textbf{гармонической}.
\Def ~
\begin{itemize}
    \item[а)]$f: G \mapsto \CC$, где $G \subseteq \CC$~--- область, называется
    \textbf{регулярной} (аналитической, голоморфной), если $\forall z \in G \
    \exists f'(z)$ (если $f$ дифференцируема в каждой точке $G$).
    \item[б)]$f: G \mapsto \CC$, $G \subseteq \CC$ называется
    \textbf{регулярной в точке $z_0 \in G$}, если $\exists r > 0: B_r(z_0)
    \subseteq G$, $f$ регулярна на $B_r(z_0)$.
\end{itemize}
\Def
Функция, заданная на произвольном множестве, \textbf{регулярна} на нем, если она
регулярна в каждой его точке.
\Example
\begin{align*}
  & w = z^n, \ n \geq 2, \ n \in \NN
\end{align*}
Производная существует в любой точке:
\begin{align*}
  & \frac{\left( z+\Delta z \right)^n - z^n}{\Delta z} = \frac{nz^{n-1}\Delta z + \left( \Delta z \right)^2\left( P(z, \Delta z) \right)}{\Delta z} \us{\Delta z \to 0}{\to} nz^{n-1}
\end{align*}
Значит, функция всюду дифференцируема и регулярна на всей $\CC$.
\\
Аналогично можно доказать регулярность всех многочленов и рациональных функций.
\Def
\begin{align*}
  & e^z = e^x\left( \cos y + i \sin y \right), \ z = x+iy
\end{align*}
Функция периодична с периодом $2 i \pi$; ее компоненты $u = e^x \cos y$, $v =
e^x \sin y$.
Выполняются УКР:
\begin{align*}
  & \begin{cases}
      u_x = e^x \cos y = v_y \\
      v_x = e^x \sin y = -u_y
  \end{cases}
  & % \left\{ \begin{matrix}
    %       u_x = e^x \cos y = v_y \\
    %       v_x = e^x \sin y = -u_y
    %   \end{matrix} \right.
\end{align*}
Значит, $f'(z) = u_x + iv_x = u+iv = e^z$ (аналогично $(e^{iz})' = ie^{iz}$,
$(e^{-iz})' = -ie^{iz}$).
\\
Эти функции всюду дифференцируемы, а значит, и регулярны.
\Def
Тригонометрические функции:
\begin{align*}
  & \cos z = \frac{e^{iz} + e^{-iz}}{2}, \ \sin z = \frac{e^{iz} -e^{-iz}}{2i}
\end{align*}
Для них выполняются известные свойства тригонометрических функций:
\begin{align*}
  & \cos^2 z + \sin^2 z = 1
\end{align*}
\begin{align*}
  & \sin(z_1+z_2) = \sin z_1 \cos z_2 + \sin z_2 \cos z_1, \ \cos(z_1+z_2) = \cos z_1 \cos z_2 - \sin z_1 \sin z_2
\end{align*}
\begin{align*}
  & \left( \sin z \right)' = \cos z, \ \left( \cos z \right)' = - \sin z
\end{align*}
Но ограниченность не выполняется.
\\
Раскроем функцию через $x$ и $y$.
\begin{align*}
  & \sin z = \sin(x+iy) = \sin x \cos iy + \cos x \sin iy = \sin x \ch y + i \cos x \sh y \\
  & \Downarrow \\
  & u(x,y) = \sin x \ch y, \ v = \cos x \sh y
\end{align*}
Тогда
\begin{align*}
  & \left| \sin z \right| = \sqrt{\sin^2 x \ch^2 y + \cos^2 x \sh^2 y} = \sqrt{(1-\cos^2 x)\ch^2 y + (\ch^2 y - 1)\cos^2 x} = \\
  & = \sqrt{\ch^2 y - \cos^2 x} \us{y \to \infty}{\sim} \ch x \us{y \to \infty}{\to} \infty
\end{align*}
\theorem
Пусть $f: G \mapsto \CC$ регулярна в $G$, $f = u + iv$, $u,v \in C^2(G)$. Тогда
$u$ и $v$ являются гармоническими.
\pr
\begin{align*}
  & \Delta u = \frac{\partial}{\partial x}\left( \frac{\partial u}{\partial x} \right)+ \frac{\partial}{\partial y} \left( \frac{\partial u}{\partial y} \right) = \frac{\partial^2 v}{\partial x \partial y} - \frac{\partial^2 v}{\partial y \partial x} = 0
\end{align*}
Переход объясняется условием Коши-Римана для регулярной $f$.
\\
Аналогично можно доказать гармоничность $v$.
\\
Пара функций, связанных УКР, называется \textbf{сопряженными гармоническими}.
\theorem
Пусть $G$~--- односвязная область, $u: G \mapsto \RR$ гармоническая. Тогда
сущетвует регулярная $f$, такая, что $\Real f(z) = u(x,y)$.
\pr
Пусть
\begin{align*}
  & P(x,y) = -\frac{\partial u}{\partial y}, \ Q(x,y) = \frac{\partial u}{\partial x}
\end{align*}
Хотим найти $v$, такую, что
\begin{align*}
  & \begin{cases}
      \dst \frac{\partial v}{\partial x} = P(x,y) \\
      \dst \frac{\partial v}{\partial y} = Q(x,y) 
      \end{cases}
\end{align*}
(из УКР на $f$).
\begin{align*}
  & dv = Pdx+Qdy
\end{align*}
Для односвязной $G$ и потенциального векторного поля $(P,Q)$, т.~е. такого, что
выполняется условие:
\begin{align*}
  & \frac{\partial P}{\partial x} -\frac{\partial Q}{\partial x}' = - \frac{\partial^2 u}{\partial y^2} - \frac{\partial^2 u}{\partial x^2} = 0
\end{align*}
В таком случае
\begin{align*}
  & v(x,y) = \int_{(x_0,y_0)}^{(x,y)} P dx + Q dy + const
\end{align*}
Полученая функция $v$ гармоническая. тогда $f(z) = u+iv$ будет регулярной в
каждой точке $G$.
\\
Аналогично можно построить регулярную функцию $f$ по данной гармонической $v$.
\Example
Пусть
\begin{align*}
  & u = xy, \ \Delta u = 0
\end{align*}
Значит, $u$ гармоническая. УКР:
\begin{align*}
  & v_x = -u_y = -x \Rightarrow v(x,y) = \frac{-x^2}{2} + \varphi(y)
\end{align*}
\begin{align*}
  & v_y = u_x = y \Rightarrow v(x,y) = \frac{-x^2}{2} + \frac{y^2}{2} + C
\end{align*}
Значит,
\begin{align*}
  & \varphi'(y) = y = v_y
\end{align*}
Тогда
\begin{align*}
  & v(x,y) = \frac{y^2-x^2}{2} + C
\end{align*}
Значит,
\begin{align*}
  & f(x,y)= xy+i\left( \frac{y^2-x^2}{2} + C \right) = -\frac{i}{2} \left( x^2-y^2+2ixy \right) + iC = -\frac{iz}{2} + iC
\end{align*}

\newpage
\section{Целочисленная арифметика}

\subsection{Уровни абстракции <<железа>>}
\begin{itemize}
    \item <<p>> и <<n>> полупроводники, проводники, конденсаторы и резисторы
    \item Транзисторы и диоды
    \item Логические элементы
    \item Арифметическо-логические юниты
    \item Центральный процессор
    \item Компьютерная система
\end{itemize}

Компьютер состоит из элементарных блоков, из которых строятся электрические
элементы. Они в свою очередь образуют логические элементы.

Ключевая часть компьютера - центральный процессор. А главная часть ЦП - арифметическо-логический блок, который выполняет все вычисления.

\subsection{Типы сигналов}

\begin{Def}
    \underline{Сигналы} --- это изменения напряжения в электрической сети. Их можно регистрировать
и дальше обрабатывать.
\end{Def}

Сигналы бывают двух видов:
\begin{itemize}
    \item \textbf{Аналоговый.} Сила сигнала изменяется непрерывно с изменением напряжения.
    \item \textbf{Цифровой.} Есть всего два состояния - высокий, либо низкий.
\end{itemize}


\subsection{Дискретные сигналы}

\begin{itemize}
    \item \textbf{Логический <<0>>.} Все, что ниже определенного значения
        напряжения --- это нулевой уровень.
    \item \textbf{Логическая <<1>>.}
    \item \textbf{Неопределенное состояние.} Напряжение между 0 и 1, не приводящее к изменению состояния.
    \item \textbf{Hi-Z состояние.} Физически какой-то контакт внутри процессора отключается от общей схемы, чтобы не оказывать влияние на все остальные компоненты.
\end{itemize}

\subsection{Арифметико-логический юнит}

Базовый юнит центрального процессора для реализации операций над числами:
\begin{itemize}
	\item Битовая логика
	\item Побитовые сдвиги
	\item Сложение и вычитание
	\item Умножение и деление
\end{itemize}

 Управление ходом программы --- тоже арифметические операции. 
 
 \subsection{Представление чисел и переполнения}
 Было на семинарах...
 
 \subsection{Значения из нескольких байтов}
 
 \begin{itemize}
 	\item Минимальный адресуемый обьем данных --- 8 бит.
 	\item За один такт процессора он может адресовать одно \textbf{машинное слово} процессора - 32 или 64 бита.
 \end{itemize}

\subsection{Big-Endian и Little-Endian}

\begin{itemize}
	\item \textbf{Big-Endian:}
	Выглядит как человеческое представление байтов.
	сначала старшие байты, потом - младшие
	Пример: большинство RISC процессоров по умолчанию.
	\item \textbf{Little-Endian:}
	сначала младшие байты, потом старшие
	Пример: Процессоры Intel и ARM.
\end{itemize}

\subsection{Ковертирование endian-ов}

Беззнаковые типы можно конвертировать просто побитовыми сдвигами. Достаточно переставить байты в обратном порядке.

Со знаковыми типами есть проблемы, потому что там первый бит отвечает за знак.

\subsection{Упаковка структур}

Размер любой структуры должен быть кратек размеру машинного слова. Это связано с тем оптимизацией времени чтения/записи. 

Поэтому если в структурке хранится два поля - на 16 бит и на 8 бит, то стуктура будет занимать 32 бита, а не 16 + 8 = 24. Ведь размер машинного слова в x86 равен 32 битам.

Если же прописать у структуры специальный атрибут, компилятор уложит поля вплотную.

\subsection{Где используется атрибут packed?}

\begin{itemize}
	\item Бинарные данные
	\item Сетевые пакеты
	\item Драйвера
\end{itemize}

В остальных случаях использование packed ненужно и даже вредно.

\newpage
\section{Команды процессора}

\subsection{Процессор}

Процессор умеет только манипулировать с целыми числами. Значение адреса - тоже число.

\subsection{Процессор с точки зрения пользователя}

У процессора есть:
\begin{itemize}
	\item Регистры - eax, ebx, ...
	\item Флаги - ZF, CF, SF
	\item Указатель на текущую команду (PC - pointer count)
\end{itemize}

\subsection{Команды x86}

nop - 0x00
halt - 0x10
call Dest - 0x80, Dest - 5 байтов
...

Команды кодируются переменным количеством байт. Количество байт может доходить до 6.

\subsection{Представление структур программы в виде простых инструкций}

Было на семинарах, на ассембелере писать уже должны уметь...

\subsection{CISC}

\begin{Def}
	\underline{CISC} --- \textbf{C}omplex \textbf{I}nstruction \textbf{S}et \textbf{C}omputing
\end{Def}

Процессоры: x86, Z80, PDP-11/VAX

Это тип архитектуры процессора, где:
\begin{itemize}
	\item Команд много, на все случаи жизни
	\item Кодируются переменным количеством байт
	\item Разные режимы адресации
	\item Упрощают жизнь программисту на ассемблере
\end{itemize}

Например, инструкция \textbf{loop} Address позволяет реализовать цикл одной командой.

\subsection{Конвейер}

\begin{itemize}
	\item Загрузка инструкции
	\item Декодирование команды
	\item Выполнение
	\item Доступ к памяти
	\item Записать результат в регистр
\end{itemize}

\subsection{Проблемы ковейеризации}

\begin{itemize}
	\item Инструкции имеют разную длину в байтах
	\item Разные инструкции выполняются за разное число тактов
	\item Работа с памятью и с регистрами задействует разные блоки процессора
\end{itemize}

Одна из задач компилятора с точки зрения оптимизации - минимизировать время простоя при конвейеризации. Для этого компилятор может менять порядок команд, если это не изменит результата, но поможет с ковейеризацией.

\subsection{RISC}

\begin{Def}
	\underline{RISC} --- \textbf{R}educed \textbf{I}nstruction \textbf{S}et \textbf{C}omputing
\end{Def}

Процессоры: ARM, AVR, MIPS, PowerPC

\begin{itemize}
	\item Только простейшие инструкции
	\item Инструкции фиксированной длины
	\item Почти у всех инструкций одинаковое время выполнения
	\item Адресация только по регистрам. Не можем читать напрямую из памяти
\end{itemize}

\subsection{Современные RICS процессоры}

\begin{itemize}
	\item \textbf{PowerPC:} Playstation, XBox, суперкомпьютеры IBM
	\item \textbf{MIPS:} WiFi/Bluetooth-адаптеры, процессоры в телевизорах
	\item \textbf{ARM:} Смартфоны
	\item \textbf{AVR:} Arduino
\end{itemize}

\subsection{Вещественная арифметика, IEEE 754}

Было на семинарах. Очень очень лень это все дублировать...





\newpage
\section{Additional information}

\end{document}